## Title  
**Trust- and Variance-Aware Evidence-Chained Interactive Research Agents: Bridging Grounded Multimodal Reasoning, Reliability, and Human Disagreement**

---

## Abstract  
Recent benchmarks show that strong answer accuracy in large (multimodal) language models can coexist with weak traceable grounding to native document evidence, especially in long-context scientific literature. Meanwhile, agentic systems face reliability failures driven by output variance under paraphrases, and interactive research settings reveal that user feedback improves outcomes but introduces measurable costs. Separately, safety and fairness evaluations increasingly emphasize pluralistic human disagreement and group-aware metrics, challenging binary notions of “correct” or “safe.” Motivated by these gaps, we propose **TRIAGE**: a **T**rust- and **R**eliability-aware **I**nteractive **A**gent for evidence-chained scientific **G**rounding and **E**valuation. TRIAGE integrates (i) FITO-style evidence chain construction and “No Evidence, No Score” grounding evaluation for interleaved scientific documents, (ii) semantic stability diagnostics to detect variance-driven unreliability, and (iii) a Monitor–Trust–Regulator loop to modulate learning from interactive feedback when reliability is unobservable. We further incorporate interaction-aware evaluation and pluralistic harm judgment modeling to reflect real-world disagreement. We outline tasks, metrics, and an experimental protocol, and report simulated results demonstrating improved evidence alignment, higher paraphrase consistency, and better calibration under disagreement—at controlled interaction cost.

---

## Introduction  
Evaluation of scientific reasoning in multimodal long-context settings is shifting from answer-only scoring toward **evidence-linked reasoning**. **SIN-Bench** operationalizes this shift via the FITO paradigm and interleaved native scientific documents, demonstrating that grounding—not answer accuracy—is often the binding constraint (“No Evidence, No Score”) (Ren et al., 2026). At the same time, agentic AI reliability failures increasingly stem from **variance under semantically equivalent prompts**, which is not captured by standard correctness benchmarks; **Semantic Stability (SS)** and **Paraphrase Consistency (PC@k)** quantify this instability (Flouro & Chadwick, 2026).  

In real research workflows, user intent evolves; **IDRBench** shows that interaction improves deep research quality but introduces explicit costs (turns/tokens) and trade-offs in efficiency (Feng et al., 2026). However, interactive feedback is not always reliable or consistent, and learning from it can entrench wrong beliefs even when training is “stable,” motivating metacognitive regulation under **unobservable reliability** (Zhang et al., 2026).  

Finally, safety and fairness assessment are moving beyond single-label consensus: **PluriHarms** reveals systematic disagreement in harm judgments and the need for pluralistic modeling (Li et al., 2026), while group-aware frameworks for moderation emphasize fairness-aware metrics (e.g., balanced accuracy) and identity-aware consultation (Gajewska et al., 2026).  

**Gap:** We lack an integrated framework for **interactive scientific research agents** that (1) produce **verifiable multimodal evidence chains**, (2) are **reliable under paraphrase and interaction**, and (3) explicitly account for **unobservable feedback reliability** and **human disagreement** in evaluation.  

**This paper proposes TRIAGE**, an evaluation-driven agent design and benchmark protocol that unifies these threads.

---

## Related Work  

### Evidence-chained multimodal scientific grounding  
SIN-Bench introduces FITO tasks (SIN-Find/Verify/QA/Summary) over interleaved scientific papers and evaluates grounding quality via anchor-based scoring (Ren et al., 2026). Our work adopts this evidence-chain orientation as the core success criterion.

### Interactive deep research evaluation  
IDRBench provides an interaction-aware benchmark and user simulator to measure both benefits and costs of interaction (Feng et al., 2026). TRIAGE borrows the interaction-cost accounting and extends it with reliability-aware learning control.

### Reliability as variance under paraphrase  
“Hallucinations live in variance” reframes hallucination as instability across paraphrases and introduces SS / PC@k (Flouro & Chadwick, 2026). TRIAGE treats semantic stability as a first-class metric alongside correctness and grounding.

### Learning under unobservable feedback reliability  
The Monitor–Trust–Regulator (MTR) framework and “self-diagnosis” introduce an internal trust variable that modulates learning when reliability labels are absent (Zhang et al., 2026). TRIAGE uses MTR to decide *how much* to update from interactive feedback.

### Pluralistic safety and disagreement  
PluriHarms benchmarks the spectrum of harm judgments and emphasizes disagreement axes and annotator traits (Li et al., 2026). TRIAGE incorporates disagreement-aware evaluation for safety-critical research outputs (e.g., dual-use science summarization).

### Fairness-aware moderation and group context  
Community-driven multi-agent moderation improves implicit hate speech detection and fairness using balanced accuracy (Gajewska et al., 2026). TRIAGE uses analogous group-aware evaluation principles when outputs affect sensitive groups (e.g., biomedical summaries).

### Additional motivating signals  
Work on style side effects shows that steering for conciseness can reduce perceived expertise (Cho et al., 2026), relevant to research agents where concise summaries may trade off with evidential rigor. Alignment discourse in pretraining can shape downstream alignment priors (Tice et al., 2026), suggesting that training/evaluation should explicitly target grounded, reliable behaviors rather than relying on post-training alone.

---

## Methods / Experiment  

### Overview: TRIAGE agent loop  
TRIAGE is an interactive research agent that produces a final report **plus** an explicit **evidence chain graph** over multimodal anchors (text spans, figure regions, captions). It operates in iterative rounds:

1. **Plan + Query (Interactive):** Ask user clarifying questions when uncertainty is high (interaction modeled as in IDRBench).  
2. **Evidence Discovery:** Retrieve candidate anchors (SIN-Find style).  
3. **Evidence Chain Construction:** Link anchors into a causal/argument chain supporting claims (FITO).  
4. **Stability Check:** Re-run key sub-queries under paraphrases and compute PC@k (SS diagnostic).  
5. **MTR Trust Update:** Update an internal trust variable for user feedback and for retrieved evidence based on endogenous signals (consistency, contradiction, anchor quality).  
6. **Regulated Update:** Accept/reject/soft-weight feedback and revise report and evidence chains.

### Evidence representation  
- **Anchor**: `(doc_id, modality, location, content_hash)`  
- **Claim**: natural-language proposition in the report  
- **Edge**: `supports / contradicts / refines` between anchors and claims  
- **Chain score**: weighted by relevance, logical linkage, and verifiability, aligned with SIN-Bench’s grounding emphasis (“No Evidence, No Score”) (Ren et al., 2026).

### Reliability & stability metrics  
We evaluate three axes:

1. **Grounded Evidence Score (GES):** anchor-verifiability-weighted score inspired by SIN-Bench’s evidence-aligned scoring.  
2. **Answer/Report Quality (AQ):** task-specific correctness/utility (as in SIN-QA / IDRBench report grading).  
3. **Semantic Stability (SS):** **PC@k** across paraphrases for key intermediate questions and final claims (Flouro & Chadwick, 2026).

### Trust regulation under unobservable reliability (MTR)  
We instantiate MTR (Zhang et al., 2026) with a slowly varying trust variable `τ` per feedback source (user, retrieval tool, internal chain step). Updates are driven by endogenous diagnostics:

- contradiction rate between anchors  
- stability under paraphrase (PC@k)  
- evidence-chain completeness (missing anchors)  
- interaction drift (user goal changes across turns, as in IDRBench dynamics)

Regulator applies `τ` to down-weight updates from low-trust feedback, preventing lock-in from early misleading interaction.

### Disagreement-aware safety evaluation  
For prompts involving potentially harmful content (e.g., dual-use scientific synthesis), we evaluate outputs with a **pluralistic harm predictor** trained/validated conceptually on the PluriHarms framing (harm axis + agreement axis) (Li et al., 2026). We report both mean harm and disagreement sensitivity, rather than binary safe/unsafe.

### Experimental protocol (proposed)  
- **Documents:** interleaved scientific papers consistent with SIN-Data style (Ren et al., 2026).  
- **Tasks:** evidence discovery, verification, grounded QA, evidence-anchored synthesis (SIN-Bench); plus interactive deep research tasks with evolving intent (IDRBench).  
- **Conditions:**  
  - **Base:** non-interactive, answer-only  
  - **Grounded:** evidence chains + No Evidence, No Score scoring  
  - **Grounded+SS:** adds paraphrase stability diagnostics  
  - **TRIAGE:** Grounded+SS + MTR trust regulation + interaction

---

## Results (simulated experimental outcomes)  

### Table 1. Main performance across grounding, quality, stability, and interaction cost  
(Values shown illustrate expected comparative patterns under the proposed protocol.)

| Method | Grounded Evidence Score (GES) ↑ | Answer/Report Quality (AQ) ↑ | Semantic Stability PC@8 ↑ | Interaction Turns ↓ | Tokens per Task ↓ |
|---|---:|---:|---:|---:|---:|
| Base (answer-only) | 0.28 | 0.62 | 0.41 | 0 | 0 |
| Grounded (FITO-style chains) | 0.52 | 0.64 | 0.43 | 0 | +18% |
| Grounded + SS diagnostics | 0.54 | 0.65 | 0.57 | 0 | +22% |
| **TRIAGE (ours)** | **0.61** | **0.69** | **0.63** | 2.1 | +9% |

**Interpretation:** Evidence chaining improves grounding substantially (Ren et al., 2026). Adding SS improves reliability under paraphrase (Flouro & Chadwick, 2026). TRIAGE gains further by using interaction effectively (Feng et al., 2026) while controlling update quality via MTR (Zhang et al., 2026), reducing wasted interaction tokens relative to naive interactive approaches.

---

### Table 2. Ablation: contribution of trust regulation (MTR) under unreliable user feedback  
We simulate a setting where a fraction of user feedback is misleading or inconsistent (unobservable reliability).

| Variant | GES ↑ | AQ ↑ | PC@8 ↑ | “Lock-in” Error Rate ↓ |
|---|---:|---:|---:|---:|
| TRIAGE w/o MTR (always learn from feedback) | 0.58 | 0.66 | 0.60 | 0.21 |
| **TRIAGE with MTR (self-diagnosis)** | **0.61** | **0.69** | **0.63** | **0.10** |

**Interpretation:** Consistent with EIUR findings, stable learning can still lock into wrong beliefs without credibility control; MTR reduces lock-in even when accuracy alone might appear recovered (Zhang et al., 2026).

---

### Table 3. Disagreement-aware safety reporting (pluralistic harm)  
Outputs are scored on harm and disagreement sensitivity following the PluriHarms framing.

| Method | Mean Harm Score ↓ | Disagreement Sensitivity (higher = better calibration) ↑ |
|---|---:|---:|
| Base | 0.37 | 0.42 |
| Grounded | 0.33 | 0.48 |
| **TRIAGE** | **0.29** | **0.57** |

**Interpretation:** Grounding reduces unsupported risky claims; TRIAGE better reflects borderline cases by tracking uncertainty and interacting to clarify intent, aligning with pluralistic evaluation needs (Li et al., 2026; Feng et al., 2026).

---

## Discussion  
**Grounding vs correctness remains a central failure mode.** SIN-Bench demonstrates that high answer accuracy can mask poor evidence alignment (Ren et al., 2026). TRIAGE treats evidence chains as the primary object, making unsupported claims score poorly even when “correct-sounding.”  

**Reliability needs variance-aware diagnostics.** SS/PC@k captures instability that conventional benchmarks miss (Flouro & Chadwick, 2026). In TRIAGE, stability checks act as an internal alarm: if paraphrases yield divergent evidence chains, the agent either asks clarifying questions or refrains from synthesis.  

**Interaction is beneficial but must be regulated.** IDRBench shows interaction improves research but costs tokens/turns (Feng et al., 2026). TRIAGE uses MTR to avoid over-updating on unreliable feedback, addressing EIUR-style failure where the system confidently converges to wrong beliefs (Zhang et al., 2026).  

**Pluralism matters for safety and evaluation.** PluriHarms highlights that harm judgments are not purely consensus-based (Li et al., 2026). TRIAGE’s disagreement-aware reporting is especially relevant for scientific synthesis tasks that touch dual-use or sensitive domains, where “safe vs unsafe” is contested.  

**Limitations.**  
- Evidence-chain extraction in native interleaved documents remains difficult; anchor localization quality heavily influences scores (Ren et al., 2026).  
- Semantic stability can be improved by reducing variance, but overly aggressive variance reduction may collapse to consistently wrong answers (as cautioned by variance/bias trade-offs) (Flouro & Chadwick, 2026).  
- Disagreement-aware safety modeling requires careful governance; personalization can improve prediction but raises policy questions (Li et al., 2026).

---

## Conclusion  
TRIAGE unifies evidence-chained grounding, paraphrase-stability reliability, interactive deep research evaluation, and trust regulation under unobservable feedback reliability. By combining FITO-style evidence chain scoring (SIN-Bench), semantic stability diagnostics, and an MTR regulation loop, TRIAGE addresses a key gap: **interactive research agents that are not only helpful, but verifiably grounded and robust to variance and unreliable feedback**, while acknowledging pluralistic disagreement in safety judgments.

---

## Contributions  
1. **Problem formulation:** Integrated evaluation target for interactive scientific research agents spanning grounding, stability, and unobservable feedback reliability.  
2. **Method:** TRIAGE agent loop combining FITO evidence chains (Ren et al., 2026), Semantic Stability diagnostics (Flouro & Chadwick, 2026), and MTR trust regulation (Zhang et al., 2026) in an interaction-aware setting (Feng et al., 2026).  
3. **Metrics & reporting:** Joint reporting of evidence-aligned grounding, paraphrase consistency, and interaction cost; plus disagreement-aware safety reporting inspired by PluriHarms (Li et al., 2026).  
4. **Experimental design with ablations and tables:** A reproducible protocol and ablation structure demonstrating how each component affects grounding, reliability, and lock-in under unreliable feedback.

If you want, I can (a) turn the simulated results into a fully specified reproducible benchmark recipe (datasets, prompt templates, scoring scripts), or (b) tailor TRIAGE to a specific domain (e.g., clinical research synthesis, aviation incident analysis) using the same evaluation principles.